\section{Derivation of the Complete-Data Likelihood}
\label{app:complete_data_likelihood}

This appendix derives the complete-data likelihood for the different
distributions considered for the latent number of competing causes. The
derivation is first presented in a general form and is then specialized to
the Bernoulli, Binomial, Negative Binomial, and Poisson distributions.

\subsection{Generic Complete-Data Likelihood}

Recall that, for the $i$th individual, $N_i$ denotes the latent number of
competing causes and
\begin{equation}
T_i^*=\min\{Z_{i1},\ldots,Z_{iN_i}\},
\end{equation}
where, conditional on $N_i=n_i$, the latent event times
$Z_{i1},\ldots,Z_{in_i}$ are assumed to be independent and identically
distributed with proper survival function $S(t)$ and density $f(t)$.

Let $C_i$ denote the censoring time, and define the observed time as
\begin{equation}
T_i=\min\{T_i^*,C_i\}.
\end{equation}
The censoring indicator is denoted by $\delta_i$, where
$\delta_i=1$ indicates that an event is observed at $t_i$, whereas
$\delta_i=0$ indicates that the observation is right censored at $t_i$.

First, conditional on $N_i=n_i$, the survival probability of the latent
failure time $T_i^*$ is
\begin{equation}
P(T_i^*>t_i\mid N_i=n_i)
=
P(Z_{i1}>t_i,\ldots,Z_{in_i}>t_i)
=
S(t_i)^{n_i}.
\label{eq:conditional_survival_N}
\end{equation}

When $\delta_i=1$, an event is observed at $t_i$, so that
$T_i^*=t_i$. Since $T_i^*$ is the minimum of the $n_i$ latent event times,
its conditional density is obtained by differentiating its survival
function:
\begin{align}
f(t_i\mid N_i=n_i)
&=
-\frac{d}{dt_i}S(t_i)^{n_i} \nonumber\\
&=
n_i f(t_i)S(t_i)^{n_i-1}.
\label{eq:conditional_density_N}
\end{align}

When $\delta_i=0$, the observation is right censored at $t_i$, which implies
that $T_i^*>t_i$. Therefore, the corresponding conditional contribution
is the survival probability given in Equation~\eqref{eq:conditional_survival_N}.

The complete data include the latent realization $N_i=n_i$. Thus, defining
\begin{equation}
    \mathcal{D}_{\mathrm{comp}} = \left\{ (t_i,\delta_i,n_i) \right\}^{n}_{i=1}\text{,}
\end{equation}
the contribution of the $i$th individual to the complete-data likelihood
can be written as
\begin{equation}
\label{eq:joint_complete_factorization}
L_{c,i}
=
P(N_i=n_i)\,
P(T_i,\delta_i\mid N_i=n_i).
\end{equation}

Combining the two possible values of $\delta_i$ gives
\begin{equation}
\label{eq:generic_complete_contribution}
L_{c,i}
=
P(N_i=n_i)
\left[
n_i f(t_i)S(t_i)^{n_i-1}
\right]^{\delta_i}
\left[
S(t_i)^{n_i}
\right]^{1-\delta_i}.
\end{equation}

The first factor in Equation~\eqref{eq:generic_complete_contribution},
$P(N_i=n_i)$, corresponds to the probability mass associated with the
latent number of competing causes. The remaining factors describe the
contribution of the observed event or censoring time conditional on this
latent value.

Combining the powers of $S(t_i)$ gives
\begin{align}
L_{c,i}
&=
P(N_i=n_i)
n_i^{\delta_i}
f(t_i)^{\delta_i}
S(t_i)^{(n_i-1)\delta_i+n_i(1-\delta_i)}
\nonumber\\
&=
P(N_i=n_i)
n_i^{\delta_i}
f(t_i)^{\delta_i}
S(t_i)^{n_i-\delta_i}.
\label{eq:generic_complete_contribution_simplified}
\end{align}

Assuming independence across individuals, the complete-data likelihood is
therefore
\begin{equation}
\label{eq:generic_complete_likelihood}
L_c
\left(
\boldsymbol{\eta};
\mathcal{D}_{\mathrm{comp}}
\right)
=
\prod_{i=1}^n
P(N_i=n_i)
n_i^{\delta_i}
f(t_i)^{\delta_i}
S(t_i)^{n_i-\delta_i},
\end{equation}
where $\boldsymbol{\eta}$ denotes the collection of model parameters.

Consequently, the complete-data log-likelihood is
\begin{equation}
\label{eq:generic_complete_loglikelihood}
\ell_c
\left(
\boldsymbol{\eta};
\mathcal{D}_{\mathrm{comp}}
\right)
=
\sum_{i=1}^n
\left[
\log P(N_i=n_i)
+
\delta_i\log n_i
+
\delta_i\log f(t_i)
+
(n_i-\delta_i)\log S(t_i)
\right].
\end{equation}

The expression in Equation~\eqref{eq:generic_complete_loglikelihood}
provides the common structure for all distributions of the latent number
of competing causes considered in this work. The distribution-specific
complete-data log-likelihoods are obtained by substituting the corresponding
probability mass function $P(N_i=n_i)$ into this expression.

\subsection{Bernoulli Distribution}

Suppose that
\begin{equation}
N_i\sim\operatorname{Bernoulli}(\theta_i),
\qquad 0<\theta_i<1,
\end{equation}
so that
\begin{equation}
P(N_i=n_i)
=
\theta_i^{n_i}(1-\theta_i)^{1-n_i},
\qquad n_i\in\{0,1\}.
\label{eq:bern_complete_pmf}
\end{equation}

Substituting Equation~\eqref{eq:bern_complete_pmf} into the generic
complete-data likelihood in Equation~\eqref{eq:generic_complete_likelihood}
gives
\begin{align}
L_c
\left(
\boldsymbol{\eta};
\mathcal{D}_{\mathrm{comp}}
\right)
&=
\prod_{i=1}^n
\theta_i^{n_i}
(1-\theta_i)^{1-n_i}
n_i^{\delta_i}
f(t_i)^{\delta_i}
S(t_i)^{n_i-\delta_i}.
\label{eq:bern_complete_likelihood}
\end{align}

Taking logarithms yields the complete-data log-likelihood
\begin{align}
\ell_c
\left(
\boldsymbol{\eta};
\mathcal{D}_{\mathrm{comp}}
\right)
&=
\sum_{i=1}^n
\Big[
n_i\log\theta_i
+
(1-n_i)\log(1-\theta_i)
+
\delta_i\log n_i
\nonumber\\
&\qquad\qquad
+
\delta_i\log f(t_i)
+
(n_i-\delta_i)\log S(t_i)
\Big].
\label{eq:bern_complete_loglikelihood}
\end{align}

Notice that when $\delta_i=1$, an event is observed and therefore
$n_i\geq1$. In the Bernoulli case, this implies $n_i=1$. Hence, the term
$\delta_i\log n_i$ is well-defined for every observed event. When
$\delta_i=0$, this term vanishes and its value does not need to be
evaluated.

\subsection{Binomial Distribution}

Now suppose that the latent number of competing causes follows a
Binomial distribution, $N_i\sim\operatorname{Binomial}(K,\theta_i^*)\text{, } 0<\theta_i^*<1$,
where $K$ is a positive integer. Its probability mass function is
\begin{equation}
P(N_i=n_i)
=
{K\choose n_i}
(\theta_i^*)^{n_i}
(1-\theta_i^*)^{K-n_i},
\qquad
n_i=0,1,\ldots,K.
\label{eq:binomial_complete_pmf}
\end{equation}

Substituting Equation~\eqref{eq:binomial_complete_pmf} into the generic
complete-data likelihood in Equation~\eqref{eq:generic_complete_likelihood}
gives
\begin{align}
L_c
\left(
\boldsymbol{\eta};
\mathcal{D}_{\mathrm{comp}}
\right)
&=
\prod_{i=1}^n
{K\choose n_i}
(\theta_i^*)^{n_i}
(1-\theta_i^*)^{K-n_i}
n_i^{\delta_i}
f(t_i)^{\delta_i}
S(t_i)^{n_i-\delta_i}.
\label{eq:binomial_complete_likelihood}
\end{align}

Taking logarithms yields the complete-data log-likelihood
\begin{align}
\ell_c
\left(
\boldsymbol{\eta};
\mathcal{D}_{\mathrm{comp}}
\right)
&=
\sum_{i=1}^n
\Bigg[
\log {K\choose n_i}
+
n_i\log\theta_i^*
+
(K-n_i)\log(1-\theta_i^*)
\nonumber\\
&\qquad\qquad
+
\delta_i\log n_i
+
\delta_i\log f(t_i)
+
(n_i-\delta_i)\log S(t_i)
\Bigg].
\label{eq:binomial_complete_loglikelihood}
\end{align}

As in the Bernoulli case, when $\delta_i=1$, an event is observed and
therefore $n_i\geq1$. Consequently, the term $\delta_i\log n_i$ is
well-defined for all observations for which an event is observed. When
$\delta_i=0$, this term vanishes and does not need to be evaluated.

\subsection{Negative Binomial Distribution}

Consider the Negative Binomial parameterization adopted in this work,
with $N_i\sim\operatorname{NB}(\alpha,\theta_i)$,
whose probability mass function is
\begin{equation}
P(N_i=n_i)
=
\frac{
\Gamma(\alpha^{-1}+n_i)
}{
\Gamma(\alpha^{-1})\,n_i!
}
\left(
\frac{\alpha\theta_i}
{1+\alpha\theta_i}
\right)^{n_i}
(1+\alpha\theta_i)^{-1/\alpha},
\qquad
n_i=0,1,2,\ldots,
\label{eq:nb_complete_pmf}
\end{equation}
under the parameter restrictions required for this parameterization.

Substituting Equation~\eqref{eq:nb_complete_pmf} into the generic
complete-data likelihood in Equation~\eqref{eq:generic_complete_likelihood}
gives
\begin{align}
L_c
\left(
\boldsymbol{\eta};
\mathcal{D}_{\mathrm{comp}}
\right)
&=
\prod_{i=1}^n
\Bigg[
\frac{
\Gamma(\alpha^{-1}+n_i)
}{
\Gamma(\alpha^{-1})\,n_i!
}
\left(
\frac{\alpha\theta_i}
{1+\alpha\theta_i}
\right)^{n_i}
(1+\alpha\theta_i)^{-1/\alpha}
\nonumber\\
&\qquad\qquad\qquad
\times
n_i^{\delta_i}
f(t_i)^{\delta_i}
S(t_i)^{n_i-\delta_i}
\Bigg].
\label{eq:nb_complete_likelihood}
\end{align}

Taking logarithms and separating the terms gives the complete-data
log-likelihood
\begin{align}
\ell_c
\left(
\boldsymbol{\eta};
\mathcal{D}_{\mathrm{comp}}
\right)
&=
\sum_{i=1}^n
\Bigg[
\log\Gamma(\alpha^{-1}+n_i)
-
\log\Gamma(\alpha^{-1})
-
\log(n_i!)
\nonumber\\
&\qquad
+
n_i
\log\left(
\frac{\alpha\theta_i}
{1+\alpha\theta_i}
\right)
-
\frac{1}{\alpha}
\log(1+\alpha\theta_i)
\nonumber\\
&\qquad
+
\delta_i\log n_i
+
\delta_i\log f(t_i)
+
(n_i-\delta_i)\log S(t_i)
\Bigg].
\label{eq:nb_complete_loglikelihood}
\end{align}

As in the Bernoulli and Binomial cases, when $\delta_i=1$, an event is
observed and therefore $n_i\geq1$. Hence, the term
$\delta_i\log n_i$ is well-defined whenever it contributes to the
log-likelihood. When $\delta_i=0$, this term vanishes and does not need
to be evaluated.

\subsection{Poisson Distribution}

Finally, suppose that the latent number of competing causes follows a
Poisson distribution, $N_i\sim\operatorname{Poisson}(\theta_i)\text{, } \theta_i>0$.
The corresponding probability mass function is
\begin{equation}
P(N_i=n_i)
=
\frac{\theta_i^{n_i}e^{-\theta_i}}{n_i!},
\qquad
n_i=0,1,2,\ldots.
\label{eq:poisson_complete_pmf}
\end{equation}

Substituting Equation~\eqref{eq:poisson_complete_pmf} into the generic
complete-data likelihood in Equation~\eqref{eq:generic_complete_likelihood}
gives
\begin{align}
L_c
\left(
\boldsymbol{\eta};
\mathcal{D}_{\mathrm{comp}}
\right)
&=
\prod_{i=1}^n
\frac{\theta_i^{n_i}e^{-\theta_i}}{n_i!}
n_i^{\delta_i}
f(t_i)^{\delta_i}
S(t_i)^{n_i-\delta_i}.
\label{eq:poisson_complete_likelihood}
\end{align}

Taking logarithms yields the complete-data log-likelihood
\begin{align}
\ell_c
\left(
\boldsymbol{\eta};
\mathcal{D}_{\mathrm{comp}}
\right)
&=
\sum_{i=1}^n
\Big[
n_i\log\theta_i
-
\log(n_i!)
-
\theta_i
+
\delta_i\log n_i
\nonumber\\
&\qquad\qquad
+
\delta_i\log f(t_i)
+
(n_i-\delta_i)\log S(t_i)
\Big].
\label{eq:poisson_complete_loglikelihood}
\end{align}

As in the previous cases, when $\delta_i=1$, an event is observed and
therefore $n_i\geq1$, so that the term $\delta_i\log n_i$ is well-defined.
When $\delta_i=0$, this term vanishes and does not need to be evaluated.

Thus, the complete-data log-likelihoods for the four distributions
considered in this work share the same contribution from the latent
competing event times. They differ only through the probability mass
function specified for the latent number of competing causes $N_i$.
